\chapter{Implementation}
\label{ch:implementation}

\section{Pipeline Overview}
\label{sec:impl-pipeline}

Figure~\ref{fig:pipeline} shows the final pipeline. It differs from the pipeline
described in the original proposal in exactly the ways described in
Chapter~\ref{ch:methodology}: there is no DPO fine-tuning stage, no MoDPO branch,
and no cross-summary sentence-alignment step. The pipeline has four stages. First,
a fixed number of CNN/DailyMail articles are selected deterministically
(Section~\ref{sec:impl-data}). Second, two candidate summaries are generated per
article at different sampling temperatures (Section~\ref{sec:impl-generation}).
Third, the same candidate pool is scored twice, once holistically and once at the
sentence level with GCA aggregation, to produce two separate preference datasets
(Section~\ref{sec:impl-preferences}). Fourth, a Bradley--Terry reward model is
trained independently on each preference dataset and the two are compared
(Section~\ref{sec:impl-rm-training}).

\begin{figure}[htbp]
\centering
\begin{tikzpicture}[
  node distance=7mm and 8mm,
  box/.style={draw, rounded corners, align=center, font=\footnotesize,
    minimum width=34mm, minimum height=9mm, inner sep=2mm},
  shared/.style={box, fill=gray!12},
  hol/.style={box, fill=blue!8},
  gca/.style={box, fill=orange!10},
  end/.style={box, fill=gray!12, minimum width=48mm},
  arr/.style={-{Stealth[length=2mm]}, thick},
]

% shared stages
\node[shared] (data)   {1. Data Loader\\\footnotesize CNN/DailyMail, seeded subset};
\node[shared, below=of data]  (gen)   {2. Candidate Generation\\\footnotesize Mistral-7B, $T{=}0.7$ / $T{=}1.0$};
\node[shared, below=of gen]   (pair)  {3. Candidate Pairing\\\footnotesize Summary A, Summary B};

% branch anchor row
\node[hol, below left=12mm and -8mm of pair]  (hscore) {4A. Holistic Scoring\\\footnotesize AlignScore on full summary};
\node[gca, below right=12mm and -8mm of pair] (seg)    {4B. Sentence Segmentation};

\node[hol, below=of hscore] (hpref)  {5A. Holistic Preference\\\footnotesize margin $=0$};
\node[gca, below=of seg]    (gscore) {5B. Sentence-Level Scoring\\\footnotesize AlignScore per sentence};

\node[gca, below=of gscore] (gagg)   {6B. GCA Aggregation\\\footnotesize $\alpha{=}0.0$ (simple mean)};
\node[gca, below=of gagg]   (gpref)  {7B. GCA Preference\\\footnotesize margin $=0$};

\node[hol, below=18mm of hpref] (hrm) {8A. Bradley--Terry RM\\\footnotesize \emph{RM-Holistic}, 5-fold CV};
\node[gca, below=of gpref]      (grm) {8B. Bradley--Terry RM\\\footnotesize \emph{RM-GCA}, 5-fold CV};

\node[end, below=12mm of grm] (eval) {9. Evaluation and Comparison\\\footnotesize pairwise accuracy, bootstrap CI, Wilcoxon test};

% arrows
\draw[arr] (data) -- (gen);
\draw[arr] (gen) -- (pair);
\draw[arr] (pair) -| (hscore);
\draw[arr] (pair) -| (seg);
\draw[arr] (hscore) -- (hpref);
\draw[arr] (seg) -- (gscore);
\draw[arr] (gscore) -- (gagg);
\draw[arr] (gagg) -- (gpref);
\draw[arr] (hpref) -- (hrm);
\draw[arr] (gpref) -- (grm);
\draw[arr] (hrm) |- (eval);
\draw[arr] (grm) -- (eval);

\end{tikzpicture}
\caption[The final experimental pipeline]{The final pipeline. Shared setup (grey) produces one candidate pool from
which two independent preference-construction branches run: holistic (blue) and
GCA (orange). Each branch trains its own Bradley--Terry reward model under an
identical recipe, and the two are compared in the final evaluation step. Compare
to the pipeline described in the original proposal (Chapter~\ref{ch:methodology}),
which additionally included a DPO fine-tuning stage and a cross-summary
sentence-alignment step, neither of which appears here.}
\label{fig:pipeline}
\end{figure}

\section{Data}
\label{sec:impl-data}

All experiments use CNN/DailyMail \citep{hermann2015teaching}, version \texttt{3.0.0}
as distributed through the Hugging Face \texttt{datasets} library. Subset selection
(\url{src/data/subset.py}, function \texttt{select\_subset}) filters the chosen
split to articles under 8{,}000 characters and reference summaries under 2{,}000
characters, then draws a fixed-size sample from the eligible pool using Python's
\texttt{random.Random} seeded deterministically. Each sample is assigned an ID of
the form \texttt{cnn\_\{index\}\_\{hash\}}, where \texttt{hash} is the first 12
hexadecimal characters of the SHA-256 digest of the article text
(\texttt{make\_sample\_id}). This ties every sample's identifier to its content
rather than to its position in the dataset, so the same article always receives the
same ID regardless of which subset size or seed selected it.

Table~\ref{tab:impl-subsets} lists every subset used across the project. The
$n=1{,}000$, $n=5{,}000$, and $n=10{,}000$ subsets all use subset seed 200 and the
same filtering parameters, differing only in the number of samples requested. They
are therefore \emph{nested rather than disjoint}: the smaller sets are almost
entirely contained in the larger ones. Chapter~\ref{ch:methodology},
Section~\ref{sec:methodology-nesting} quantifies the overlap and states its
methodological consequence. Earlier subsets at seeds 42 and 100 were drawn with
different seeds and are not part of the final comparison.

\begin{table}[htbp]
\centering
\caption[CNN/DailyMail subsets used across the project]{CNN/DailyMail subsets used across the project.}
\label{tab:impl-subsets}
\begin{tabularx}{\textwidth}{lclX}
\toprule
$n$ & Subset seed & Purpose & Notes \\
\midrule
200    & 42  & Early DPO-based prototype & Superseded; see
  Chapter~\ref{ch:results}, Section~\ref{sec:results-early} \\
495    & 100 & First Bradley--Terry RM training set & Drawn with a different seed
  from the $n=200$ set \\
1{,}000 & 200 & Final repeated-run campaign & 6 RM training runs; seeds 42, 42, 7, 100,
  314, 2026 \\
5{,}000 & 200 & Scale-robustness check & Single training run \\
10{,}000 & 200 & Scale-robustness check & Single training run \\
\bottomrule
\end{tabularx}
\end{table}

\section{Candidate Generation}
\label{sec:impl-generation}

Candidate summaries are generated with Mistral-7B-Instruct-v0.3, loaded in
\texttt{bfloat16} without quantization (\url{src/generation/candidates.py}). On the
A100-SXM4-40GB GPUs used for this project, the model occupies roughly 14\,GB in
\texttt{bfloat16}, comfortably within budget, so 4-bit quantization, though
supported in the code via \texttt{BitsAndBytesConfig}, is disabled in the
configuration actually used (\texttt{load\_in\_4bit: false} in
\url{configs/generation.yaml}). The tokenizer uses left-padding, required for
batched generation with a decoder-only model, and articles are truncated to 1{,}024
input tokens.

Each article is summarized twice, using the fixed prompt template
\texttt{"Summarize the following news article in a concise paragraph.\ \ldots"},
with two different sampling configurations: a low-temperature setting
($T=0.7$, top-$p=0.9$) and a high-temperature setting ($T=1.0$, top-$p=0.95$),
generating up to 256 new tokens each. These two outputs form the candidate pair
(``summary A'' and ``summary B'') that both the holistic and GCA preference
construction stages score. Generation is resumable: the script checks which
sample IDs already appear in the output file and only generates candidates for the
remainder, which matters in practice on a shared cluster where a job may need to be
resubmitted after a queue timeout.

\section{Preference Construction}
\label{sec:impl-preferences}

Both preference-construction regimes are implemented in
\url{src/judging/build_reward_preferences.py} and share a single judge object,
\texttt{RewardModelJudge} (\url{src/judging/reward_model_judge.py}), configured to
use the AlignScore backend.

In the holistic regime (\texttt{judge\_holistic\_pair}), each full candidate
summary is scored once against the article with \texttt{judge.score(article,\
summary)}, and the two resulting scores are compared with a margin guard
(\texttt{make\_decision}): summary A is preferred if its score exceeds summary B's
by more than the margin, summary B is preferred in the opposite case, and the pair
is marked \texttt{no\_preference} otherwise. With the final margin of 0
(Chapter~\ref{ch:methodology}, Section~\ref{sec:methodology-evolution}), every
pair with a nonzero score
difference produces a decision.

In the GCA regime (\texttt{judge\_gca\_pair}), each candidate summary is first
split into sentences with a regex-based splitter, then every sentence is scored
independently against the article with \texttt{judge.score\_sentences(article,\
sentences)}. The resulting list of per-sentence scores is passed to
\texttt{aggregate\_sentence\_scores(scores,\ alpha)}
(Chapter~\ref{ch:background}, Equation~\ref{eq:gca-general}) to produce a single
GCA score per summary, and the same margin-guarded comparison determines the
preference. Every output record retains the full sentence-level detail, including
each sentence's individual AlignScore score and a count of ``low-faithfulness''
sentences (score below 0.5), which supports the qualitative disagreement analysis
in Chapter~\ref{ch:results}.

The locked final configuration is invoked as:
\begin{quote}
\small\ttfamily
python src/judging/build\_reward\_preferences.py\\
\hspace*{1em}--mode both --alpha 0.0 --margin 0\\
\hspace*{1em}--alignscore-evaluation-mode nli\\
\hspace*{1em}--alignscore-backbone FacebookAI/roberta-base
\end{quote}

\section{Reward-Model Training}
\label{sec:impl-rm-training}

Reward-model training is implemented in \url{src/reward_model/train.py}
(\texttt{BradleyTerryRewardModel}, \texttt{bradley\_terry\_loss}) and orchestrated
by \url{src/reward_model/run_training.py}, which trains the holistic and GCA
conditions from a single invocation. All results reported in
Chapter~\ref{ch:results} use $k$-fold cross-validation
(\texttt{\_kfold\_cv} in \texttt{run\_training.py}): the preference pairs are
shuffled using a seeded \texttt{random.Random} instance, split into $k=5$
contiguous folds, and for each fold a freshly initialized reward model is trained
on the other four folds and evaluated on the held-out one. One implementation
detail is worth stating precisely, since it affects how ``training seed'' should
be interpreted. The \texttt{--seed} argument is passed to Python's
\texttt{random} module, which governs the fold-assignment shuffle. PyTorch's own
generator is not seeded anywhere in the cross-validation path
(\texttt{\_kfold\_cv}), so every source of stochasticity drawn from it is
uncontrolled: the randomly initialized reward head, the shuffling of training
batches by the \texttt{DataLoader}, and the dropout masks applied during
training. Changing \texttt{--seed} therefore changes the fold assignment, while
these other sources vary between process executions regardless of the seed value.
Together they explain why two executions launched with the same seed value
produced different results (Chapter~\ref{ch:results},
Section~\ref{sec:results-final-campaign}), and they are discussed as a source of
run-to-run variation in Chapter~\ref{ch:discussion},
Section~\ref{sec:disc-stochasticity}.

The final training configuration, used for every reported reward-model result, is:
\url{FacebookAI/roberta-base} backbone, 5 epochs, batch size 8, learning rate
$2\times10^{-5}$ with a linear warmup over 10\% of total steps, maximum sequence
length 512 tokens (article truncated to 2{,}000 characters before tokenization,
concatenated with the summary), and 5-fold cross-validation. This is confirmed by
\url{slurm/train_rm_1000.sh}, the Slurm script used for the headline
$n=1{,}000$ result. One inconsistency in the codebase is worth flagging explicitly,
since it could otherwise look like an error in the reported results: the in-code
default for \texttt{--backbone} in both \texttt{train.py} and
\texttt{run\_training.py} is \url{microsoft/deberta-v3-base}, a leftover from an
earlier version of the training script. Every Slurm script used to produce a
result reported in this thesis overrides this default and passes
\texttt{--backbone FacebookAI/roberta-base} explicitly; the \texttt{deberta-v3-base}
default was never actually used to produce a reported number. The one exception is
the very first Bradley--Terry result on the $n=495$ set
(Chapter~\ref{ch:results}, Section~\ref{sec:results-rm-baseline}), which predates
the standardization on \texttt{roberta-base} and is reported as such.

\paragraph{Source-document truncation.} The reward model does not see the full
source article. Two limits apply in sequence. The article string is first
truncated to its leading 2{,}000 characters
(\texttt{--max-article-chars}, applied in \texttt{PreferencePairDataset}), and the
truncated article is then concatenated with the candidate summary and tokenized
under a 512-token cap (\texttt{--max-length}), so the article prefix is shortened
further to make room for the summary. Subset selection admits articles up to
8{,}000 characters (Section~\ref{sec:impl-data}), so for the longer articles in
the pool the reward model observes well under half of the source text.

This matters for interpreting the reward-model results. The evaluator that
produced the preference labels scores summaries against the article under its own
chunking scheme, whereas the reward model must reconstruct those labels from a
prefix only. Where the evidence supporting or contradicting a claim lies beyond
the retained prefix, that evidence is simply unavailable to the reward model, and
the corresponding preference is not recoverable from its input in principle rather
than merely in practice. How large this effect is was not measured here, and the
near-chance accuracies reported in Chapter~\ref{ch:results} should not be
attributed to truncation without evidence; it is one plausible contributing factor
among several discussed in Chapter~\ref{ch:discussion},
Section~\ref{sec:disc-weak-signal}. Varying \texttt{--max-article-chars} while
holding everything else fixed would isolate it, and is recorded as future work.

\section{Compute Infrastructure}
\label{sec:impl-infra}

All experiments run on MOGON NHR (\texttt{mogon-nhr-01.zdv.uni-mainz.de}), the
University of Mainz's high-performance computing cluster, using Slurm partition
\texttt{a100dl} under account \texttt{nhr-haloed}, on NVIDIA A100-SXM4-40GB GPUs.
All Slurm scripts producing results reported in Chapter~\ref{ch:results} exclude
one cluster node found to have a faulty CUDA runtime, via
\texttt{\#SBATCH --exclude}; this affects scheduling only and not the computation
performed.

\section{Reproducibility}
\label{sec:impl-repro}

The pipeline is config-driven throughout: dataset, generation, and judging
parameters live in \url{configs/*.yaml} and can be overridden per run from the
command line, and every script accepts explicit flags for the parameters that
matter for a given experiment (margin, $\alpha$, AlignScore evaluation mode,
reward-model hyperparameters, and so on) rather than hardcoding them.

Seeding is partial, and the limits should be stated exactly. Subset selection and
cross-validation fold assignment are explicitly seeded through Python's
\texttt{random} module. PyTorch's generator is not seeded in the
cross-validation path, so weight initialization, batch ordering and dropout are
not fixed, and exact run-level reproduction from the command-line seed alone is
not guaranteed: two executions with identical arguments can differ. Every
pipeline script writes a JSON run-metadata file
(\url{src/utils/logging.py}, \texttt{save\_run\_metadata}) recording its
configuration, output artifact paths, and a timestamped run ID, so that any
reported number can in principle be traced back to the exact invocation that
produced it.

The configuration required to reproduce the $n=1{,}000$ result is given in full
in Table~\ref{tab:methodology-summary} (Chapter~\ref{ch:methodology}), together
with the six training seeds listed in Chapter~\ref{ch:results},
Table~\ref{tab:seed-validation}, and is implemented by
\url{slurm/train_rm_1000.sh} and \url{slurm/mode_nli_seed_confirm.sh}. The
codebase, including all scripts referenced in this chapter, is available at the
project's GitHub repository: \url{https://github.com/hasnat23/master-thesis-rlaif-gca}.
